\section{One simultaneous reversal for two projectivities}
\begin{theorem}\label{E.SimultaneousReversal}
Let $F$ be a field and let $A,B\in\mathrm{GL}_2(F)$ have no common
invariant line over an algebraic closure. Then $X=AB-BA$ is invertible,
$X^2$ is scalar, and conjugation by $[X]$ simultaneously inverts
$[A],[B]\in\mathrm{PGL}_2(F)$. In particular, every generating pair of
$\mathrm{PSL}_2(q)$ or $\mathrm{PGL}_2(q)$ admits an automorphism of
order dividing two which inverts both generators.
\end{theorem}
\begin{proof}
The Cayley--Hamilton identity gives
\[
 XA+AX=\operatorname{tr}(A)X,\qquad
 XB+BX=\operatorname{tr}(B)X.
\]
Indeed, the first left side is $A^2B-BA^2$, and the second is
$AB^2-B^2A$. Thus $XA=\operatorname{adj}(A)X$ and
$XB=\operatorname{adj}(B)X$. These identities hold in every
characteristic. Also $\operatorname{tr}(X)=0$, whence
$X^2=-\det(X)I$.

To see that $X$ is invertible, extend scalars to an algebraic closure.
Neither $A$ nor $B$ is scalar, since otherwise they share an eigenline.
If $A=\operatorname{diag}(a,d)$ with $a\ne d$ and
$B=\left(\begin{smallmatrix}u&v\\w&z\end{smallmatrix}\right)$,
then $\det(AB-BA)=(a-d)^2vw$. Its vanishing makes either coordinate
line common and invariant. Otherwise $A$ has one Jordan block and,
after rescaling its off-diagonal entry, the commutator is
$\left(\begin{smallmatrix}w&z-u\\0&-w\end{smallmatrix}\right)$.
Singularity again makes the unique eigenline of $A$ invariant under
$B$. Both possibilities contradict the hypothesis. Since
$\operatorname{adj}(A)=\det(A)A^{-1}$, the two matrix identities now
prove the projective assertion.

For the finite projective groups, the upper and lower elementary
unipotents have distinct unique fixed points over the algebraic closure.
Both belong to $\mathrm{PSL}_2(q)$. Hence a pair generating either
projective group cannot share an invariant line. Conjugation by $[X]$
normalizes $\mathrm{PGL}_2(q)$ and its determinant-square subgroup
$\mathrm{PSL}_2(q)$, and its square is the identity.
\end{proof}

For rank-three rotations $a,b$, apply the theorem to the generating pair
$a,ab$. Since $ab$ is an involution, the resulting involutory automorphism
sends $a$ to $a^{-1}$ and fixes $ab$, hence sends $b$ to $a^2b$.
This is exactly the rank-three rotation reversal rule. The retained
rotation-group correspondence therefore excludes a chiral polyhedron
with either projective group as its full automorphism group.
The correspondence is the retained geometric input; the preceding theorem
supplies the complete algebraic reversal. It does not require the product
of the two generators to be an involution or their cyclic subgroups to
have trivial intersection. This recovers Elliott--Leemans--O'Reilly-Regueiro,
Theorem~4.1, by one argument, including the small fields.

For clarity, the source's auxiliary involution count requires a
characteristic qualification. For $u:x\mapsto x+1$, an inverting
projectivity fixes infinity and has the form $x\mapsto ax+b$.
The conjugation equation forces $a=-1$. In odd characteristic all $q$
choices of $b$ give nonidentity involutions. In characteristic two these
are translations; the choice $b=0$ is the identity and the other $q-1$
choices are involutions. Thus the literal count in its Lemma~4.3 is $q-1$
when the characteristic is two. This correction does not affect the
principal nonexistence theorem.
